% Generated by tools/render.py -- do not edit.
\usevoicingcolor{ink}
\coverpage{\coversubject{Bass}{fretyellow}}{}
\usevoicingcolor{fretyellow}
\sectiondivider{Bass}{(B) E A D G}
\usevoicingcolor{fretyellow}
\usecofsize{7}
\begin{circlepage}{Bass}
\begin{circleoffifths}{Bass}
  \cofsegment{0}{C}{B1 (A3)}{Am}{E5}{$\natural$}{up}
  \cofsegment{1}{G}{E3}{Em}{B5 (E0)}{1\#}{up}
  \cofsegment{2}{D}{B3 (A5)}{Bm}{B0 (E7)}{2\#}{up}
  \cofsegment{3}{A}{E5}{F\#m}{B7 (E2)}{3\#}{up}
  \cofsegment{4}{E}{B5 (E0)}{C\#m}{B2 (A4)}{4\#}{down}
  \cofsegment{5}{B}{B0 (E7)}{G\#m}{E4}{5\#}{down}
  \cofsegment{6}{Gb}{B7 (E2)}{Ebm}{B4 (A6)}{6$\flat$}{down}
  \cofsegment{7}{Db}{B2 (A4)}{Bbm}{E6}{5$\flat$}{down}
  \cofsegment{8}{Ab}{E4}{Fm}{B6 (E1)}{4$\flat$}{down}
  \cofsegment{9}{Eb}{B4 (A6)}{Cm}{B1 (A3)}{3$\flat$}{down}
  \cofsegment{10}{Bb}{E6}{Gm}{E3}{2$\flat$}{up}
  \cofsegment{11}{F}{B6 (E1)}{Dm}{B3 (A5)}{1$\flat$}{up}
\end{circleoffifths}
\circlefootnote{String and fret for each root: \frets{A3} is the third fret of the A string.\\ Lowest that falls under the hand.\\ Five-string first, four-string in parentheses.}
\end{circlepage}
\begin{bookpage}{Root Positions}
\pagesubtitle{Bass}
\begin{rootmap}{5}{7pt}
\rootmaphead{{\bfseries (B)} & {\bfseries E} & {\bfseries A} & {\bfseries D} & {\bfseries G}}
\rootmaprow{C}{\frets{1} & \frets{8} & \frets{3} & \frets{10} & \frets{5}}
\rootmaprow{Db}{\frets{2} & \frets{9} & \frets{4} & \frets{11} & \frets{6}}
\rootmaprow{D}{\frets{3} & \frets{10} & \frets{5} & \frets{0} & \frets{7}}
\rootmaprow{Eb}{\frets{4} & \frets{11} & \frets{6} & \frets{1} & \frets{8}}
\rootmaprow{E}{\frets{5} & \frets{0} & \frets{7} & \frets{2} & \frets{9}}
\rootmaprow{F}{\frets{6} & \frets{1} & \frets{8} & \frets{3} & \frets{10}}
\rootmaprow{Gb}{\frets{7} & \frets{2} & \frets{9} & \frets{4} & \frets{11}}
\rootmaprow{G}{\frets{8} & \frets{3} & \frets{10} & \frets{5} & \frets{0}}
\rootmaprow{Ab}{\frets{9} & \frets{4} & \frets{11} & \frets{6} & \frets{1}}
\rootmaprow{A}{\frets{10} & \frets{5} & \frets{0} & \frets{7} & \frets{2}}
\rootmaprow{Bb}{\frets{11} & \frets{6} & \frets{1} & \frets{8} & \frets{3}}
\rootmaprow{B}{\frets{0} & \frets{7} & \frets{2} & \frets{9} & \frets{4}}
\end{rootmap}
\begin{rootmapnote}
Fret for the root of each key on each string.\\
Everything else is measured from there.
\end{rootmapnote}
\end{bookpage}
\begin{bookpage}{What to Play Under}
\pagesubtitle{Bass \quad degrees from the root}
\begin{degreetable}
\degreerow{major}{\frets{R} \frets{3} \frets{5}}
\degreerow{m}{\frets{R} \frets{b3} \frets{5}}
\degreerow{5}{\frets{R} \frets{5}}
\degreerow{+}{\frets{R} \frets{3} \frets{\#5}}
\degreerow{\degsign{}}{\frets{R} \frets{b3} \frets{b5}}
\degreerow{sus2}{\frets{R} \frets{2} \frets{5}}
\degreerow{sus4}{\frets{R} \frets{4} \frets{5}}
  \chordgap
\degreerowfirst{6}{\frets{R} \frets{3} \frets{5} \frets{6}}
\degreerow{m6}{\frets{R} \frets{b3} \frets{5} \frets{6}}
\degreerow{6/9}{\frets{R} \frets{9} \frets{3} \frets{5} \frets{6}}
  \chordgap
\degreerowfirst{7}{\frets{R} \frets{3} \frets{5} \frets{b7}}
\degreerow{maj7}{\frets{R} \frets{3} \frets{5} \frets{7}}
\degreerow{m7}{\frets{R} \frets{b3} \frets{5} \frets{b7}}
\degreerow{\degsign{}7}{\frets{R} \frets{b3} \frets{b5} \frets{6}}
\degreerow{m7b5}{\frets{R} \frets{b3} \frets{b5} \frets{b7}}
\degreerow{7sus4}{\frets{R} \frets{4} \frets{5} \frets{b7}}
\degreerow{7b5}{\frets{R} \frets{3} \frets{b5} \frets{b7}}
\degreerow{7\#5}{\frets{R} \frets{3} \frets{\#5} \frets{b7}}
  \chordgap
\degreerowfirst{9}{\frets{R} \frets{9} \frets{3} \frets{5} \frets{b7}}
\degreerow{maj9}{\frets{R} \frets{9} \frets{3} \frets{5} \frets{7}}
\degreerow{m9}{\frets{R} \frets{9} \frets{b3} \frets{5} \frets{b7}}
\degreerow{add9}{\frets{R} \frets{9} \frets{3} \frets{5}}
\degreerow{madd9}{\frets{R} \frets{9} \frets{b3} \frets{5}}
\degreerow{7b9}{\frets{R} \frets{b9} \frets{3} \frets{5} \frets{b7}}
\degreerow{7\#9}{\frets{R} \frets{b3} \frets{3} \frets{5} \frets{b7}}
\degreerow{m11}{\frets{R} \frets{9} \frets{b3} \frets{4} \frets{5} \frets{b7}}
\end{degreetable}
\begin{rootmapnote}
Play the root, and the fifth if there is one.\\
Add the tone that names the chord, the \frets{b3} or the \frets{b7}, only when it wants hearing.\\
The rest belongs to whoever is playing chords.
\end{rootmapnote}
\end{bookpage}
\begin{bookpage}{Patterns}
\pagesubtitle{Bass \quad counted from the root}
\begin{patternlist}
\patternitem{Root + fifth + octave}{\begin{tabular}{@{}r@{\hskip 2mm}l@{\hskip 3mm}l@{\hskip 3mm}l@{}}\patternrow{degree} & \frets{R} & \frets{5} & \frets{8}\\
\patternrow{string} & \patterncell{R} & \patterncell{1+R} & \patterncell{2+R}\\
\patternrow{fret} & \patterncell{R0} & \patterncell{R2} & \patterncell{R2}
\end{tabular}}{The default. Works under anything in the key.}
\patternitem{Major arpeggio}{\begin{tabular}{@{}r@{\hskip 2mm}l@{\hskip 3mm}l@{\hskip 3mm}l@{\hskip 3mm}l@{}}\patternrow{degree} & \frets{R} & \frets{5} & \frets{8} & \frets{3}\\
\patternrow{string} & \patterncell{R} & \patterncell{1+R} & \patterncell{2+R} & \patterncell{3+R}\\
\patternrow{fret} & \patterncell{R0} & \patterncell{R2} & \patterncell{R2} & \patterncell{R1}
\end{tabular}}{Major, maj7, 6.}
\patternitem{Minor arpeggio}{\begin{tabular}{@{}r@{\hskip 2mm}l@{\hskip 3mm}l@{\hskip 3mm}l@{\hskip 3mm}l@{}}\patternrow{degree} & \frets{R} & \frets{5} & \frets{8} & \frets{b3}\\
\patternrow{string} & \patterncell{R} & \patterncell{1+R} & \patterncell{2+R} & \patterncell{3+R}\\
\patternrow{fret} & \patterncell{R0} & \patterncell{R2} & \patterncell{R2} & \patterncell{R0}
\end{tabular}}{Minor, m7, m6.}
\patternitem{Dominant seventh}{\begin{tabular}{@{}r@{\hskip 2mm}l@{\hskip 3mm}l@{\hskip 3mm}l@{\hskip 3mm}l@{}}\patternrow{degree} & \frets{R} & \frets{5} & \frets{b7} & \frets{3}\\
\patternrow{string} & \patterncell{R} & \patterncell{1+R} & \patterncell{2+R} & \patterncell{3+R}\\
\patternrow{fret} & \patterncell{R0} & \patterncell{R2} & \patterncell{R0} & \patterncell{R1}
\end{tabular}}{Any 7 chord. The b7 pulls to the four.}
\patternitem{Minor seventh}{\begin{tabular}{@{}r@{\hskip 2mm}l@{\hskip 3mm}l@{\hskip 3mm}l@{\hskip 3mm}l@{}}\patternrow{degree} & \frets{R} & \frets{5} & \frets{b7} & \frets{b3}\\
\patternrow{string} & \patterncell{R} & \patterncell{1+R} & \patterncell{2+R} & \patterncell{3+R}\\
\patternrow{fret} & \patterncell{R0} & \patterncell{R2} & \patterncell{R0} & \patterncell{R0}
\end{tabular}}{The two and the six chord.}
\patternitem{Walking approach}{\begin{tabular}{@{}r@{\hskip 2mm}l@{\hskip 3mm}l@{\hskip 3mm}l@{\hskip 3mm}l@{}}\patternrow{degree} & \frets{R} & \frets{5} & \frets{6} & \frets{b7}\\
\patternrow{string} & \patterncell{R} & \patterncell{1+R} & \patterncell{1+R} & \patterncell{2+R}\\
\patternrow{fret} & \patterncell{R0} & \patterncell{R2} & \patterncell{R4} & \patterncell{R0}
\end{tabular}}{Leading into the next chord.}
\end{patternlist}
\end{bookpage}
\usevoicingcolor{ink}
\begin{bookpage}{Numbers, Major}
\pagesubtitle{Any instrument}
\begin{numberchart}
\numberhead{1}{2m}{3m}{4}{5}{6m}{7\textdegree}
\numberrow{C}{\numbercell{C} & \numbercell{Dm} & \numbercell{Em} & \numbercell{F} & \numbercell{G} & \numbercell{Am} & \numbercell{B\degsign{}}}
\numberrow{Db}{\numbercell{Db} & \numbercell{Ebm} & \numbercell{Fm} & \numbercell{Gb} & \numbercell{Ab} & \numbercell{Bbm} & \numbercell{C\degsign{}}}
\numberrow{D}{\numbercell{D} & \numbercell{Em} & \numbercell{F\#m} & \numbercell{G} & \numbercell{A} & \numbercell{Bm} & \numbercell{C\#\degsign{}}}
\numberrow{Eb}{\numbercell{Eb} & \numbercell{Fm} & \numbercell{Gm} & \numbercell{Ab} & \numbercell{Bb} & \numbercell{Cm} & \numbercell{D\degsign{}}}
\numberrow{E}{\numbercell{E} & \numbercell{F\#m} & \numbercell{G\#m} & \numbercell{A} & \numbercell{B} & \numbercell{C\#m} & \numbercell{D\#\degsign{}}}
\numberrow{F}{\numbercell{F} & \numbercell{Gm} & \numbercell{Am} & \numbercell{Bb} & \numbercell{C} & \numbercell{Dm} & \numbercell{E\degsign{}}}
\numberrow{Gb}{\numbercell{Gb} & \numbercell{Abm} & \numbercell{Bbm} & \numbercell{Cb} & \numbercell{Db} & \numbercell{Ebm} & \numbercell{F\degsign{}}}
\numberrow{G}{\numbercell{G} & \numbercell{Am} & \numbercell{Bm} & \numbercell{C} & \numbercell{D} & \numbercell{Em} & \numbercell{F\#\degsign{}}}
\numberrow{Ab}{\numbercell{Ab} & \numbercell{Bbm} & \numbercell{Cm} & \numbercell{Db} & \numbercell{Eb} & \numbercell{Fm} & \numbercell{G\degsign{}}}
\numberrow{A}{\numbercell{A} & \numbercell{Bm} & \numbercell{C\#m} & \numbercell{D} & \numbercell{E} & \numbercell{F\#m} & \numbercell{G\#\degsign{}}}
\numberrow{Bb}{\numbercell{Bb} & \numbercell{Cm} & \numbercell{Dm} & \numbercell{Eb} & \numbercell{F} & \numbercell{Gm} & \numbercell{A\degsign{}}}
\numberrow{B}{\numbercell{B} & \numbercell{C\#m} & \numbercell{D\#m} & \numbercell{E} & \numbercell{F\#} & \numbercell{G\#m} & \numbercell{A\#\degsign{}}}
\end{numberchart}
\begin{rootmapnote}
Read across: in the key on the left, the four chord is the column headed 4. Call a song in numbers and it transposes itself.
\end{rootmapnote}
\end{bookpage}
\begin{bookpage}{Numbers, Minor}
\pagesubtitle{Any instrument}
\begin{numberchart}
\numberhead{1m}{2\textdegree}{b3}{4m}{5m}{b6}{b7}
\numberrow{C}{\numbercell{Cm} & \numbercell{D\degsign{}} & \numbercell{Eb} & \numbercell{Fm} & \numbercell{Gm} & \numbercell{Ab} & \numbercell{Bb}}
\numberrow{Db}{\numbercell{Dbm} & \numbercell{Eb\degsign{}} & \numbercell{Fb} & \numbercell{Gbm} & \numbercell{Abm} & \numbercell{Bbb} & \numbercell{Cb}}
\numberrow{D}{\numbercell{Dm} & \numbercell{E\degsign{}} & \numbercell{F} & \numbercell{Gm} & \numbercell{Am} & \numbercell{Bb} & \numbercell{C}}
\numberrow{Eb}{\numbercell{Ebm} & \numbercell{F\degsign{}} & \numbercell{Gb} & \numbercell{Abm} & \numbercell{Bbm} & \numbercell{Cb} & \numbercell{Db}}
\numberrow{E}{\numbercell{Em} & \numbercell{F\#\degsign{}} & \numbercell{G} & \numbercell{Am} & \numbercell{Bm} & \numbercell{C} & \numbercell{D}}
\numberrow{F}{\numbercell{Fm} & \numbercell{G\degsign{}} & \numbercell{Ab} & \numbercell{Bbm} & \numbercell{Cm} & \numbercell{Db} & \numbercell{Eb}}
\numberrow{Gb}{\numbercell{Gbm} & \numbercell{Ab\degsign{}} & \numbercell{Bbb} & \numbercell{Cbm} & \numbercell{Dbm} & \numbercell{Ebb} & \numbercell{Fb}}
\numberrow{G}{\numbercell{Gm} & \numbercell{A\degsign{}} & \numbercell{Bb} & \numbercell{Cm} & \numbercell{Dm} & \numbercell{Eb} & \numbercell{F}}
\numberrow{Ab}{\numbercell{Abm} & \numbercell{Bb\degsign{}} & \numbercell{Cb} & \numbercell{Dbm} & \numbercell{Ebm} & \numbercell{Fb} & \numbercell{Gb}}
\numberrow{A}{\numbercell{Am} & \numbercell{B\degsign{}} & \numbercell{C} & \numbercell{Dm} & \numbercell{Em} & \numbercell{F} & \numbercell{G}}
\numberrow{Bb}{\numbercell{Bbm} & \numbercell{C\degsign{}} & \numbercell{Db} & \numbercell{Ebm} & \numbercell{Fm} & \numbercell{Gb} & \numbercell{Ab}}
\numberrow{B}{\numbercell{Bm} & \numbercell{C\#\degsign{}} & \numbercell{D} & \numbercell{Em} & \numbercell{F\#m} & \numbercell{G} & \numbercell{A}}
\end{numberchart}
\begin{rootmapnote}
The natural minor. Players often raise the seventh of the five chord to make it dominant, which turns 5m into 5.
\end{rootmapnote}
\end{bookpage}
\begin{backsheet}{Fancy Bass Chords and Their Voicings}
\tuningrow{Bass}{(B) E A D G}{The guitar's lowest four strings, an octave down.}
\end{backsheet}
